% unit: u14
% sources: sources/005-exam-1-review.txt, sources/022-exam-1-review.txt, sources/103-mit18-06scf11-ses1-14sum-pdf-download-file.txt, sources/006-exam-1-problem-solving.txt
\section{Exam 1 review}\label{sec:u14}

The first exam covers Part I, Sections~\ref{sec:u01}--\ref{sec:u13}, with the emphasis on
rectangular matrices: nullspaces, ranks, complete solutions and the four subspaces. This section
collects the essential facts in one place and then works through the review problems of the course,
which test understanding by approaching familiar ideas from unfamiliar directions, for instance by
giving the answer and asking for the matrix.

\subsection{The essentials of Part I}

\begin{itemize}
\item \textbf{Elimination} (Sections~\ref{sec:u03}--\ref{sec:u05}). Row operations are
multiplications by elementary matrices: $E_{32}E_{31}E_{21}A=U$, so $A=LU$ with the multipliers in
$L$ (with row exchanges, $PA=LU$). Gauss--Jordan reduces $\begin{bmatrix}A&I\end{bmatrix}$ to
$\begin{bmatrix}I&A^{-1}\end{bmatrix}$ when $A$ is invertible.
\item \textbf{Subspaces} (Sections~\ref{sec:u06}--\ref{sec:u07}). A subspace contains $\vb{0}$ and
is closed under addition and scalar multiplication. Intersections of subspaces are subspaces;
unions usually are not.
\item \textbf{Rank and special solutions} (Section~\ref{sec:u08}). $r=$ number of pivots;
$R=\operatorname{rref}(A)$ has the block form $\begin{bmatrix}I&F\\0&0\end{bmatrix}$ (pivot columns
first), and the special solutions are the columns of $N=\begin{bmatrix}-F\\I\end{bmatrix}$.
\item \textbf{Complete solution} (Section~\ref{sec:u09}). $A\vb{x}=\vb{b}$ is solvable iff
$\vb{b}\in\Col{A}$; then $\vb{x}=\vb{x}_p+\vb{x}_n$. Full column rank $r=n$: at most one solution;
full row rank $r=m$: at least one (Table~\ref{tab:u09-cases}).
\item \textbf{Basis and dimension} (Section~\ref{sec:u10}). Independent and spanning; every basis of
a space has the same size. Pivot columns of $A$ are a basis of $\Col{A}$.
\item \textbf{Four subspaces} (Section~\ref{sec:u11}). Dimensions $r$, $n-r$, $r$, $m-r$ for
$\Col{A}$, $\Nul{A}$, $\Col{A\T}$, $\Nul{A\T}$. Row operations preserve the row space and the
nullspace.
\item \textbf{Beyond $\R^n$} (Sections~\ref{sec:u12}--\ref{sec:u13}). Spaces of matrices and
functions; rank one means $\vb{u}\vb{v}\T$; for a connected graph, $\Nul{A}$ is the constants,
$\Nul{A\T}$ is spanned by loops, and $n-m+\ell=1$.
\end{itemize}

Figure~\ref{fig:u14-flow} is the algorithm for $A\vb{x}=\vb{b}$ that runs through all of it.

\begin{figure}[ht]\centering
\begin{tikzpicture}[>=Stealth,font=\small,node distance=0.55cm,
    box/.style={draw,rounded corners=2pt,fill=blue!5,align=center,inner sep=4pt,text width=4.2cm},
    dec/.style={draw,diamond,aspect=2.6,fill=orange!10,align=center,inner sep=1pt,text width=2.9cm}]
  \node[box] (elim) {eliminate $\begin{bmatrix}A&\vb{b}\end{bmatrix}\to\begin{bmatrix}R&\vb{d}\end{bmatrix}$;\\ count pivots: $r$};
  \node[dec,below=of elim] (test) {zero rows of $R$\\have $d_i=0$?};
  \node[box,right=1.1cm of test,text width=3.2cm,fill=gray!10] (none) {no solution:\\$\vb{b}\notin\Col{A}$};
  \node[box,below=of test] (xp) {$\vb{x}_p$: free variables $=0$,\\pivot variables from $\vb{d}$};
  \node[box,below=of xp] (xn) {special solutions $\vb{s}_1,\dots,\vb{s}_{n-r}$\\(columns of $N$)};
  \node[box,below=of xn,text width=5.4cm] (all) {$\vb{x}=\vb{x}_p+c_1\vb{s}_1+\dots+c_{n-r}\vb{s}_{n-r}$\\unique exactly when $r=n$};
  \draw[->] (elim) -- (test);
  \draw[->] (test) -- node[above] {no} (none);
  \draw[->] (test) -- node[right] {yes} (xp);
  \draw[->] (xp) -- (xn);
  \draw[->] (xn) -- (all);
\end{tikzpicture}
\caption{Solving $A\vb{x}=\vb{b}$ for any $m\times n$ matrix. When $r=m$ there are no zero rows and
the test always passes; when $r=n$ there are no special solutions.}
\label{fig:u14-flow}
\end{figure}

\subsection{Solved review problems}

\begin{example}[The dimension of a span]\label{ex:u14-span}
\emph{Nonzero vectors $\vb{u},\vb{v},\vb{w}\in\R^7$ span a subspace. What are its possible
dimensions?}
The answer is $1$, $2$ or $3$. A basis can be chosen from among the three spanning vectors, so the
dimension is at most $3$; it is at least $1$ because the vectors are nonzero. All three values
occur: $\vb{u}=\vb{v}=\vb{w}$ gives $1$, and three independent vectors give $3$.
\end{example}

\begin{example}[Stacked echelon forms]\label{ex:u14-blocks}
\emph{A $5\times3$ matrix $R$ in reduced echelon form has $r=3$ pivots. (a) Find $\Nul{R}$. (b) Find
the reduced echelon form and the rank of the $10\times3$ matrix $B=\begin{bmatrix}R\\2R\end{bmatrix}$.
(c) Do the same for the $10\times6$ matrix $C=\begin{bmatrix}R&R\\R&0\end{bmatrix}$. (d) Find
$\dim\Nul{C\T}$.}

(a) Three pivots in three columns: no free variables, so $\Nul{R}=\{\vb{0}\}$. In fact
$R=\begin{bmatrix}I_3\\0\end{bmatrix}$ is forced: three pivots equal to $1$, zeros elsewhere in
their columns, two zero rows.

(b) Subtract $2\times$(rows $1$--$5$) from rows $6$--$10$: $B\to\begin{bmatrix}R\\0\end{bmatrix}$,
which is already in reduced form. The rank is $3$, not $6$: the bottom rows of $B$ are multiples of
the top rows.

(c) Work with block rows. Subtract the top block row from the bottom one, then add the new bottom
block row to the top one, then multiply the bottom by $-1$:
\[
\begin{bmatrix}R&R\\R&0\end{bmatrix}\to\begin{bmatrix}R&R\\0&-R\end{bmatrix}\to
\begin{bmatrix}R&0\\0&-R\end{bmatrix}\to\begin{bmatrix}R&0\\0&R\end{bmatrix}.
\]
Moving the zero rows (rows $4,5$ and $9,10$) to the bottom gives
$\begin{bmatrix}I_6\\0\end{bmatrix}$. The rank is $6$.

(d) $C$ is $10\times6$ of rank $6$, so $\dim\Nul{C\T}=m-r=10-6=4$ (and $\dim\Nul{C}=6-6=0$).
\end{example}

\begin{example}[From the complete solution back to the matrix]\label{ex:u14-recover}
\emph{The complete solution of $A\vb{x}=(2,4,2)$ is}
\[
\vb{x}=\begin{bmatrix}2\\0\\0\end{bmatrix}+c\begin{bmatrix}1\\1\\0\end{bmatrix}+d\begin{bmatrix}0\\0\\1\end{bmatrix}.
\]
\emph{Find the shape of $A$, its rank, $A$ itself, and the $\vb{b}$ for which $A\vb{x}=\vb{b}$ is
solvable.}

$\vb{x}$ and $\vb{b}$ both have three components, so $A$ is $3\times3$. The two vectors multiplying
$c$ and $d$ are independent and lie in $\Nul{A}$; nothing else does (otherwise there would be more
free constants), so $\dim\Nul{A}=2$ and $r=3-2=1$. To recover $A$: $A(2,0,0)=2\,\vb{a}_1=(2,4,2)$
gives $\vb{a}_1=(1,2,1)$; $A(0,0,1)=\vb{a}_3=\vb{0}$; and $A(1,1,0)=\vb{a}_1+\vb{a}_2=\vb{0}$ gives
$\vb{a}_2=-\vb{a}_1$. Hence
\[
A=\begin{bmatrix}1&-1&0\\2&-2&0\\1&-1&0\end{bmatrix},\qquad
A\begin{bmatrix}2\\0\\0\end{bmatrix}=\begin{bmatrix}2\\4\\2\end{bmatrix}.
\]
$A\vb{x}=\vb{b}$ is solvable exactly when $\vb{b}\in\Col{A}$, i.e.\ when $\vb{b}$ is a multiple of
$(1,2,1)$. This is the case $r<m$, $r<n$: zero or infinitely many solutions.
\end{example}

\begin{example}[A factored matrix]\label{ex:u14-CD}
\emph{Let}
\[
B=CD=\begin{bmatrix}1&1&0\\0&1&0\\1&0&1\end{bmatrix}\begin{bmatrix}1&0&-1&2\\0&1&1&-1\\0&0&0&0\end{bmatrix}.
\]
\emph{Without multiplying, find a basis of $\Nul{B}$ and the complete solution of
$B\vb{x}=(1,0,1)$.}

$B$ is $3\times4$, so $\Nul{B}\subseteq\R^4$. The first factor $C$ is invertible: in a combination of
its columns giving $\vb{0}$, the second row forces the weight of column~2 to be $0$, and columns~1
and~3 are clearly independent. By Proposition~\ref{prop:u08-same-null}, $\Nul{CD}=\Nul{D}$, and $D$ is
already in reduced form with pivot columns $1,2$ and $F=\begin{bmatrix}-1&2\\1&-1\end{bmatrix}$.
The special solutions (columns of $\begin{bmatrix}-F\\I\end{bmatrix}$) are
\[
\vb{s}_1=(1,-1,1,0),\qquad\vb{s}_2=(-2,1,0,1),
\]
a basis of $\Nul{B}$ (and $\rank B=2$). For a particular solution, note that the first column of $B$
is $C$ times the first column of $D$, which is the first column of $C$, namely $(1,0,1)$, exactly the
right-hand side. So $\vb{x}_p=(1,0,0,0)$ and
\[
\vb{x}=\begin{bmatrix}1\\0\\0\\0\end{bmatrix}+c\begin{bmatrix}1\\-1\\1\\0\end{bmatrix}
+d\begin{bmatrix}-2\\1\\0\\1\end{bmatrix}.
\]
Any other solution would serve equally well as $\vb{x}_p$; this one is simply the most convenient.
\end{example}

\begin{example}[A square system with a parameter]\label{ex:u14-k}
\emph{Let $A=\begin{bmatrix}1&1&1\\1&2&3\\3&4&k\end{bmatrix}$ and $\vb{b}=(2,3,7)$. (a) For which $k$
does $A\vb{x}=\vb{b}$ have exactly one solution? (b) Infinitely many? (c) Find $A=LU$. (d) Find the
complete solution for every $k$.}

Eliminate on the augmented matrix and record the steps, since part (c) needs them. $E_{21}$ subtracts
row~1 from row~2, $E_{31}$ subtracts $3\times$row~1 from row~3, and $E_{32}$ subtracts row~2 from
row~3:
\[
\left[\begin{array}{ccc|c}1&1&1&2\\1&2&3&3\\3&4&k&7\end{array}\right]\to
\left[\begin{array}{ccc|c}1&1&1&2\\0&1&2&1\\0&1&k-3&1\end{array}\right]\to
\left[\begin{array}{ccc|c}1&1&1&2\\0&1&2&1\\0&0&k-5&0\end{array}\right].
\]
(a) The third pivot is $k-5$. For $k\ne5$, $A$ is invertible and the solution is unique.
(b) For $k=5$ the third row becomes $0=0$, consistent, and $x_3$ is free: infinitely many solutions.
(Had the last entry been nonzero, $k=5$ would give no solution at all.)
(c) $U$ is the final left block, and $L=E_{21}^{-1}E_{31}^{-1}E_{32}^{-1}$ holds the multipliers
$1,3,1$ in their positions:
\[
A=\begin{bmatrix}1&0&0\\1&1&0\\3&1&1\end{bmatrix}\begin{bmatrix}1&1&1\\0&1&2\\0&0&k-5\end{bmatrix}.
\]
(d) For $k\ne5$: $x_3=0$, then $x_2=1$, then $x_1=2-1-0=1$, so $\vb{x}=(1,1,0)$. For $k=5$: set
$x_3=c$; then $x_2=1-2c$ and $x_1=2-x_2-x_3=1+c$, so
\[
\vb{x}=\begin{bmatrix}1\\1\\0\end{bmatrix}+c\begin{bmatrix}1\\-2\\1\end{bmatrix}.
\]
Check: with $k=5$, $A(1,-2,1)=(1-2+1,\,1-4+3,\,3-8+5)=\vb{0}$.
\end{example}

\subsection{Quick checks}

The following short questions test whether the definitions are understood. Each answer is followed
by its reason.
\begin{enumerate}[1.]
\item \emph{$A$ is square and $\Nul{A}=\{\vb{0}\}$. What is $\Nul{A\T}$?} Also $\{\vb{0}\}$: $r=n=m$, so
$\dim\Nul{A\T}=m-r=0$.
\item \emph{Do the invertible $5\times5$ matrices form a subspace of the $25$-dimensional space of
$5\times5$ matrices?} No: the zero matrix is not invertible, and $I+(-I)=0$. (The singular matrices
fail too: $\operatorname{diag}(1,0)+\operatorname{diag}(0,1)=I$ in the $2\times2$ case.)
\item \emph{If $B^2=0$, must $B=0$?} No: $B=\begin{bmatrix}0&1\\0&0\end{bmatrix}$ has $B^2=0$.
\item \emph{An $n\times n$ system with independent columns is solvable for every $\vb{b}$.} True:
$r=n=m$, so $A$ is invertible and $\vb{x}=A^{-1}\vb{b}$.
\item \emph{If $m=n$, the row space equals the column space.} False: for
$B=\begin{bmatrix}0&1\\0&0\end{bmatrix}$ the row space is the line through $(0,1)$ and the column
space is the line through $(1,0)$. Only the dimensions agree. (For a symmetric matrix they are equal.)
\item \emph{$A$ and $-A$ have the same four subspaces.} True: each subspace contains $-\vb{v}$
whenever it contains $\vb{v}$.
\item \emph{If $A$ and $B$ have the same four subspaces, then $A$ is a multiple of $B$.} False: every
invertible $3\times3$ matrix has $\Col{A}=\Col{A\T}=\R^3$ and both nullspaces $\{\vb{0}\}$, yet
$I$ and $\operatorname{diag}(1,2,3)$ are not multiples of each other. What \emph{is} forced is equal
rank.
\item \emph{If two rows of $A$ are exchanged, which subspaces stay the same?} The row space and the
nullspace. The column space can change: exchanging the rows of $\begin{bmatrix}1&0\\0&0\end{bmatrix}$
moves $\Col{A}$ from the line through $(1,0)$ to the line through $(0,1)$.
\item \emph{Why can $\vb{v}=(1,2,3)$ not be both a row of $A$ and in $\Nul{A}$?} The entry of
$A\vb{v}$ in that row would be $\vb{v}\cdot\vb{v}=1+4+9=14\ne0$. In fact the row space and the
nullspace meet only in $\vb{0}$, and they are perpendicular, the first topic of Part II
(Section~\ref{sec:u16}).
\end{enumerate}

\begin{trap}
Stacking copies of rows does not add rank. In Example~\ref{ex:u14-blocks} it is tempting to say
$\rank\begin{bmatrix}R\\2R\end{bmatrix}=6$ because the matrix ``contains $R$ twice''; but rank
counts independent rows, and the rows of $2R$ are multiples of those of $R$. In
$\begin{bmatrix}R&R\\R&0\end{bmatrix}$, by contrast, the second block column breaks the dependence,
and the rank does double.
\end{trap}

\subsection{Exercises}

\begin{exercise}\label{exr:u14-dims}
$A$ is $3\times5$ of rank $3$. Give the dimensions of its four subspaces. Is $A\vb{x}=\vb{b}$
solvable for every $\vb{b}$? Can a solution be unique?
\end{exercise}
\begin{solution}
$\dim\Col{A}=3$, $\dim\Nul{A}=5-3=2$, $\dim\Col{A\T}=3$, $\dim\Nul{A\T}=3-3=0$. Since $r=m=3$,
$\Col{A}=\R^3$ and every $\vb{b}$ is reachable. Since $r<n$, there are two free variables, so every
solvable system has infinitely many solutions, never exactly one.
\end{solution}

\begin{exercise}\label{exr:u14-k}
For which numbers $k$ are $(1,1,0)$, $(0,1,1)$, $(1,0,k)$ linearly independent? For the exceptional
value, find the dependency.
\end{exercise}
\begin{solution}
Put them in the columns of $A=\begin{bmatrix}1&0&1\\1&1&0\\0&1&k\end{bmatrix}$ and eliminate:
row~2 $-$ row~1 gives $(0,1,-1)$, then row~3 $-$ row~2 gives $(0,0,k+1)$. Three pivots unless
$k=-1$, so the vectors are independent exactly when $k\ne-1$. For $k=-1$ the free variable
$x_3=1$ gives $x_2=1$ and $x_1=-1$: $-(1,1,0)+(0,1,1)+(1,0,-1)=\vb{0}$.
\end{solution}

\begin{exercise}\label{exr:u14-invertible}
$M$ is an invertible matrix. Which of the four subspaces of $A$ are unchanged in $MA$ (when $M$ is
$m\times m$)? Which are unchanged in $AM$ (when $M$ is $n\times n$)?
\end{exercise}
\begin{solution}
$MA$: the nullspace is unchanged ($MA\vb{x}=\vb{0}\iff A\vb{x}=\vb{0}$) and so is the row space
(the rows of $MA$ are combinations of the rows of $A$, and conversely via $M^{-1}$); this is
Theorem~\ref{thm:u11-row}. The column space and left nullspace may change, as row exchanges show.
$AM$: transposing, $(AM)\T=M\T A\T$ with $M\T$ invertible, so by the first part $A\T$ and $M\T A\T$
have the same nullspace and row space; that is, $AM$ has the same left nullspace and column space as
$A$. Its nullspace and row space may change. In every case the rank is unchanged.
\end{solution}
